%
\begin{isabellebody}%
\setisabellecontext{Modular{\isacharunderscore}{\kern0pt}Shifts}%
%
\isadelimtheory
%
\endisadelimtheory
%
\isatagtheory
\isakeywordONE{theory}\isamarkupfalse%
\ Modular{\isacharunderscore}{\kern0pt}Shifts\isanewline
\ \ \isakeywordTWO{imports}\ Main\isanewline
\isakeywordTWO{begin}%
\endisatagtheory
{\isafoldtheory}%
%
\isadelimtheory
%
\endisadelimtheory
%
\isadelimdocument
%
\endisadelimdocument
%
\isatagdocument
%
\isamarkupsection{Modular Arithmetic and Shift Ciphers%
}
\isamarkuptrue%
%
\endisatagdocument
{\isafolddocument}%
%
\isadelimdocument
%
\endisadelimdocument
%
\begin{isamarkuptext}%
\subsection{Introduction to Shift Encryption}
  This theory formalizes a classical modular shift cipher operating over natural numbers.
  We define the encryption transformation function and verify that sequential shifts compose additively.%
\end{isamarkuptext}\isamarkuptrue%
%
\begin{isamarkuptext}%
\paragraph{Encryption Function Definition}
  We define encryption $E(n, k, m) = (m + k) \pmod n$ where $n$ is the alphabet size, 
  $k$ is the shift key, and $m$ is the message token.%
\end{isamarkuptext}\isamarkuptrue%
\isakeywordONE{definition}\isamarkupfalse%
\ E\ {\isacharcolon}{\kern0pt}{\isacharcolon}{\kern0pt}\ {\isachardoublequoteopen}nat\ {\isasymRightarrow}\ nat\ {\isasymRightarrow}\ nat\ {\isasymRightarrow}\ nat{\isachardoublequoteclose}\ \isakeywordTWO{where}\isanewline
\ \ {\isachardoublequoteopen}E\ n\ k\ m\ {\isasymequiv}\ {\isacharparenleft}{\kern0pt}m\ {\isacharplus}{\kern0pt}\ k{\isacharparenright}{\kern0pt}\ mod\ n{\isachardoublequoteclose}%
\begin{isamarkuptext}%
\paragraph{Concrete Test Evaluation}
  We evaluate the encryption function on sample concrete inputs to verify functional behavior.%
\end{isamarkuptext}\isamarkuptrue%
\isakeywordONE{value}\isamarkupfalse%
\ {\isachardoublequoteopen}E\ {\isadigit{2}}{\isadigit{6}}\ {\isadigit{3}}\ {\isadigit{5}}{\isachardoublequoteclose}%
\begin{isamarkuptext}%
\paragraph{Composition Verification}
  We prove that double encryption under key $k_1$ followed by key $k_2$ is equivalent to 
  a single encryption under key $k_1 + k_2$.%
\end{isamarkuptext}\isamarkuptrue%
\isakeywordONE{theorem}\isamarkupfalse%
\ shift{\isacharunderscore}{\kern0pt}compose{\isacharcolon}{\kern0pt}\ {\isachardoublequoteopen}E\ n\ k{\isadigit{2}}\ {\isacharparenleft}{\kern0pt}E\ n\ k{\isadigit{1}}\ m{\isacharparenright}{\kern0pt}\ mod\ n\ {\isacharequal}{\kern0pt}\ E\ n\ {\isacharparenleft}{\kern0pt}k{\isadigit{1}}\ {\isacharplus}{\kern0pt}\ k{\isadigit{2}}{\isacharparenright}{\kern0pt}\ m\ mod\ n{\isachardoublequoteclose}\isanewline
%
\isadelimproof
\ \ %
\endisadelimproof
%
\isatagproof
\isakeywordONE{by}\isamarkupfalse%
\ {\isacharparenleft}{\kern0pt}simp\ add{\isacharcolon}{\kern0pt}\ E{\isacharunderscore}{\kern0pt}def\ mod{\isacharunderscore}{\kern0pt}simps\ algebra{\isacharunderscore}{\kern0pt}simps{\isacharparenright}{\kern0pt}%
\endisatagproof
{\isafoldproof}%
%
\isadelimproof
%
\endisadelimproof
%
\begin{isamarkuptext}%
The formal property verified by the Isabelle kernel is %
\begin{isabelle}%
\isaconst{E}\ \isavar{n}\ \isavar{k{\isadigit{2}}{\isachardot}{\kern0pt}{\isadigit{0}}}\ {\isacharparenleft}{\kern0pt}\isaconst{E}\ \isavar{n}\ \isavar{k{\isadigit{1}}{\isachardot}{\kern0pt}{\isadigit{0}}}\ \isavar{m}{\isacharparenright}{\kern0pt}\ mod\ \isavar{n}\ {\isacharequal}{\kern0pt}\ \isaconst{E}\ \isavar{n}\ {\isacharparenleft}{\kern0pt}\isavar{k{\isadigit{1}}{\isachardot}{\kern0pt}{\isadigit{0}}}\ {\isacharplus}{\kern0pt}\ \isavar{k{\isadigit{2}}{\isachardot}{\kern0pt}{\isadigit{0}}}{\isacharparenright}{\kern0pt}\ \isavar{m}\ mod\ \isavar{n}%
\end{isabelle}.%
\end{isamarkuptext}\isamarkuptrue%
%
\isadelimtheory
%
\endisadelimtheory
%
\isatagtheory
\isakeywordTWO{end}\isamarkupfalse%
%
\endisatagtheory
{\isafoldtheory}%
%
\isadelimtheory
%
\endisadelimtheory
%
\end{isabellebody}%
\endinput
%:%file=/cygdrive/d/Repositories/She Crypts/formal/week1/theories/Modular_Shifts.thy%:%
%:%10=1%:%
%:%11=1%:%
%:%12=2%:%
%:%13=3%:%
%:%27=5%:%
%:%39=8%:%
%:%40=9%:%
%:%41=10%:%
%:%45=14%:%
%:%46=15%:%
%:%47=16%:%
%:%49=18%:%
%:%50=18%:%
%:%51=19%:%
%:%53=22%:%
%:%54=23%:%
%:%56=25%:%
%:%57=25%:%
%:%59=28%:%
%:%60=29%:%
%:%61=30%:%
%:%63=32%:%
%:%64=32%:%
%:%67=33%:%
%:%71=33%:%
%:%72=33%:%
%:%81=36%:%
%:%84=36%:%
%:%92=39%:%
